% Source: https://artofproblemsolving.com/community/c3h1805277_help_with_inequality
% Stable identity: aops-topic-1805277-q01
\textbf{Minimum.} Let $s=a+b+c$. Then
\[
S+3=\sum\left(\dfrac{a}{b+c}+1\right)=s\left(\dfrac{1}{b+c}+\dfrac{1}{c+a}+\dfrac{1}{a+b}\right).
\]
By the Cauchy-Schwarz inequality,
\[
\bigl((b+c)+(c+a)+(a+b)\bigr)\left(\dfrac{1}{b+c}+\dfrac{1}{c+a}+\dfrac{1}{a+b}\right)\ge9,
\]
i.e.\ $2s\sum\dfrac{1}{b+c}\ge9$. So $S+3\ge\dfrac{9}{2}$ and $S\ge\dfrac{3}{2}$, with equality exactly when $b+c=c+a=a+b$, i.e.\ for an equilateral triangle. The minimum is $\boxed{\dfrac{3}{2}}$.

\begin{asy}
usepackage("amsmath");
size(4.5cm);
defaultpen(fontsize(10pt));
pair A=(0,0),B=(5,0),C=(1.6,3);
draw(A--B--C--cycle,linewidth(1));
label("$A$",A,SW);label("$B$",B,SE);label("$C$",C,N);
label("$c$",(A+B)/2,S);label("$a$",(B+C)/2,NE);label("$b$",(A+C)/2,NW);
label("\mbox{}",(min(currentpicture,user=true)+max(currentpicture,user=true))/2,black); // keeps the text after the figure black when LaTeX typesets the labels inline
\end{asy}

\textbf{No maximum.} By the triangle inequality $b+c>a$, so $2(b+c)>a+b+c=s$ and $\dfrac{a}{b+c}<\dfrac{2a}{s}$. Adding the three such inequalities,
\[
S<\dfrac{2a+2b+2c}{s}=2 .
\]
The bound $2$ is approached but never reached: for the triangle with sides $1,1,\varepsilon$ ($0<\varepsilon<2$),
\[
S=\dfrac{2}{1+\varepsilon}+\dfrac{\varepsilon}{2}\longrightarrow2\qquad(\varepsilon\to0).
\]
So $S$ takes values in $\left[\dfrac{3}{2},2\right)$: the minimum is $\dfrac{3}{2}$ and there is \textbf{no maximum} (the supremum $2$ is not attained).

(If $a,b,c$ were only positive numbers with $a+b>c$, the minimum would still be $\dfrac{3}{2}$, but $S$ would be unbounded: $a=t$, $b=c=1$ gives $S>\dfrac {t}{2}$.)
